\chapter{Intégration dans SAP Commerce Cloud}


\section{Introduction}

L'utilité des fonctions développées dépend de leur intégration effective dans la plateforme SAP Commerce Cloud \citeURL{sap_commerce_cloud}, où elles peuvent être exploitées directement par les utilisateurs métier et s'insérer dans les processus de gestion des produits.

Ce chapitre détaille les étapes de cette intégration. Il présente tout d'abord l'architecture de SAP Commerce Cloud afin de comprendre l'organisation de la plateforme et les mécanismes sur lesquels repose son extensibilité. Il décrit ensuite la communication entre SAP Commerce Cloud et le microservice via des services REST, les extensions apportées au modèle de données pour prendre en charge les nouveaux objets métier ainsi que les interfaces développées dans Backoffice pour permettre aux utilisateurs de piloter les opérations d'enrichissement et de validation. Enfin, il détaille l'automatisation des traitements à travers les workflows ainsi que le système de notifications et de traçabilité mis en place pour assurer le suivi des opérations.

Nous présentons ici l'automatisation des flux d'enrichissement et de validation dans le respect de l'architecture standard de la plateforme.

\section{Architecture logicielle de SAP Commerce Cloud}

Avant de présenter les mécanismes d'intégration du microservice d'intelligence artificielle, il est nécessaire de comprendre l'architecture logicielle de SAP Commerce Cloud \citeURL{sap_commerce_cloud}. L'architecture modulaire de la plateforme assure une séparation nette des responsabilités entre les différents composants. Cette organisation facilite le développement et la maintenance du système tout en offrant les points d'extension nécessaires à l'intégration de nouvelles fonctionnalités. Cette analyse structurelle permet de justifier les choix d'implémentation adoptés.
\subsection{Fondamentaux de SAP Commerce Cloud}

SAP Commerce Cloud \citeURL{sap_commerce_cloud} est une plateforme complète de gestion du commerce électronique et du commerce unifié, conçue pour orchestrer l'intégralité du cycle de vie des produits et des commandes. Au-delà de son rôle de vitrine marchande, SAP Commerce centralise la gestion du catalogue, des prix, des stocks, des promotions et l'exécution des commandes. C'est un système d'information stratégique pour le secteur du commerce de détail et de la distribution en ligne.

Comme tout système d'information complexe, SAP Commerce s'appuie sur une organisation structurée du code, que l'on désigne par le terme d'architecture. L'architecture d'une application logicielle décrit comment le système est décomposé en composants, comment ces composants interagissent et quelles responsabilités incombent à chacun. Une bonne architecture facilite la compréhension du code, la maintenance évolutive, les tests automatisés et la répartition du travail au sein d'équipes de développement.

SAP Commerce adopte une architecture en couches, un pattern architectural reconnu et très largement utilisé dans l'industrie logicielle. Le principe fondamental de cette approche est la séparation des préoccupations : chaque couche gère un aspect distinct du système, et les dépendances entre couches sont strictement contrôlées. Cette organisation rend possible la modification d'une couche sans impact sur les autres. Par exemple, remplacer la base de données relationnelle par une base NoSQL n'affecterait pas la couche métier, car elle communique avec la persistance via une abstraction (la couche d'accès aux données). De même, changer l'interface utilisateur n'impacte pas la logique métier si cette dernière est bien encapsulée.
\newpage
L'architecture en couches apporte une structure rigoureuse au projet. Premièrement, elle permet à plusieurs développeurs de travailler en parallèle sur différentes couches sans créer de conflits. Deuxièmement, elle facilite les tests : chaque couche peut être testée indépendamment des autres. Troisièmement, elle rend le système extensible : ajouter de nouvelles fonctionnalités devient une question de créer de nouveaux composants dans les couches appropriées, plutôt que de modifier du code existant. Enfin, elle favorise la réutilisabilité : une logique métier bien encapsulée dans une couche Service peut être appelée par plusieurs interfaces différentes (Backoffice, API REST, batch, etc.).

\subsection{Les couches de SAP Commerce}

L'architecture en couches de SAP Commerce Cloud est illustrée par la figure~\ref{fig:architecture_sap}. Chaque couche possède une responsabilité bien définie et communique uniquement avec les couches adjacentes, ce qui favorise la modularité, la maintenabilité et l'évolutivité de la plateforme.

\begin{figure}[H]
    \centering
    \includegraphics[width=0.61\textwidth]{architecture-sap.png}
    \caption{Architecture en couches de SAP Commerce Cloud}
    \label{fig:architecture_sap}
\end{figure}

Comme le montre la figure~\ref{fig:architecture_sap}, les requêtes émises par les utilisateurs transitent successivement par les différentes couches avant d'accéder aux données persistées. Chaque couche assure un rôle spécifique dans le traitement des opérations métier.

\subsubsection{Couche Présentation}

La couche Présentation constitue le point d'interaction entre les utilisateurs et la plateforme. Dans SAP Commerce Cloud, cette couche est principalement représentée par le Backoffice, qui permet aux utilisateurs métier d'administrer les catalogues, de gérer les produits et d'exécuter les différentes opérations disponibles. Son rôle est de présenter les informations, de recueillir les actions de l'utilisateur et de transmettre les demandes aux couches inférieures.

\subsubsection{Couche Facade}

La couche Facade agit comme un intermédiaire entre la présentation et la logique métier. Elle simplifie les échanges avec les services en proposant une interface unifiée aux composants de la couche Présentation. Cette couche permet également de centraliser les appels aux différents services nécessaires à l'exécution d'une opération.

\subsubsection{Couche Service}

La couche Service regroupe la logique métier de l'application. Elle implémente les traitements nécessaires au fonctionnement de la plateforme et coordonne les différentes opérations réalisées lors d'un processus métier. Les services sont indépendants de l'interface utilisateur, ce qui permet leur réutilisation dans différents contextes d'exécution.

\subsubsection{Couche Data Access Object (DAO)}

La couche DAO assure l'accès aux données persistées. Elle constitue une abstraction entre la logique métier et le mécanisme de persistance en encapsulant les opérations de lecture, de création, de modification et de suppression des données. Cette séparation permet d'isoler la logique métier des détails liés au stockage.

\subsubsection{Couche Modèle}

La couche Modèle représente les entités métier manipulées par SAP Commerce Cloud. Ces entités, appelées \textit{ItemTypes}, définissent la structure des données persistées ainsi que les relations existantes entre elles. Elles constituent la base sur laquelle reposent les différentes opérations réalisées par les couches supérieures.

L'organisation en couches de SAP Commerce Cloud facilite ainsi le développement de nouvelles fonctionnalités tout en préservant une séparation claire des responsabilités. Cette architecture constitue le socle sur lequel s'appuie l'intégration du microservice d'intelligence artificielle présentée dans les sections suivantes.
\section{Communication entre SAP Commerce Cloud et le module d'intelligence artificielle}

L'intégration du module d'intelligence artificielle dans SAP Commerce Cloud \citeURL{sap_commerce_cloud} repose sur une communication basée sur des API REST \citeURL{fielding_rest}. SAP Commerce exploite ainsi les fonctions d'IA externes sans compromettre la séparation des responsabilités entre les composants. La plateforme SAP Commerce reste responsable de la gestion des produits et des processus métier, tandis que le module IA réalise les traitements spécialisés liés à l'enrichissement et à la validation des données.

\subsection{Principe de communication REST}

REST (\textit{Representational State Transfer}) \citeURL{fielding_rest} est un style architectural permettant à des applications distribuées de communiquer à travers le protocole HTTP. Dans cette approche, les fonctionnalités d'un système sont exposées sous forme de ressources accessibles via des URLs appelées points d'accès ou \textit{endpoints}.

La communication repose sur des requêtes HTTP utilisant différentes méthodes selon l'opération souhaitée. Dans le cadre de ce projet, les échanges entre SAP Commerce Cloud et le module d'intelligence artificielle utilisent principalement la méthode \texttt{POST}. Cette méthode permet d'envoyer des informations vers un service distant afin de déclencher un traitement.

Une requête HTTP POST contient généralement :
\begin{itemize}
    \item l'adresse de l'API appelée (\textit{endpoint}) ;
    \item les informations nécessaires au traitement transmises dans le corps de la requête (généralement au format JSON \citeURL{json_org}) ;
    \item les paramètres techniques permettant d'identifier le format des données échangées.
\end{itemize}

Dans cette architecture, SAP Commerce joue le rôle de client en envoyant les requêtes, tandis que le module d'intelligence artificielle agit comme un serveur capable de traiter ces demandes et de retourner les résultats correspondants.

\subsection{API d'enrichissement et API de validation}

Le module d'intelligence artificielle expose deux API REST principales consommées par SAP Commerce Cloud. Ces interfaces permettent d'intégrer les fonctionnalités d'enrichissement automatique et de contrôle qualité directement dans le cycle de gestion des produits.

\subsubsection{API d'enrichissement}

L'API d'enrichissement permet à SAP Commerce de transmettre les informations nécessaires à la génération automatique des contenus produit. Les données envoyées au module d'intelligence artificielle regroupent les informations descriptives et techniques nécessaires au traitement, telles que les informations générales du produit, sa marque, ses catégories, ses caractéristiques ainsi que les champs à enrichir et les langues cibles.

À partir de ces éléments, le module IA génère les contenus demandés et effectue les traitements associés, notamment la génération automatique des descriptions et leur adaptation aux différentes langues sélectionnées. Les résultats retournés sont ensuite exploités par SAP Commerce afin de mettre à jour les informations produit concernées.

\subsubsection{API de validation}

L'API de validation permet d'évaluer automatiquement la qualité des données produit existantes. Elle reçoit les informations nécessaires aux contrôles définis par le moteur de validation, notamment les données descriptives, les caractéristiques techniques et les éléments soumis aux différentes règles d'analyse.

Après traitement, cette API retourne les résultats de validation comprenant l'évaluation globale de la qualité des données, les anomalies détectées, leur niveau de sévérité ainsi que les corrections ou recommandations associées. Ces informations permettent ensuite à SAP Commerce d'assurer le suivi de la qualité des produits et d'appliquer les actions nécessaires.

\subsection{Flux global de communication}

Le fonctionnement global des échanges est illustré dans la figure~\ref{fig:communication_ia}. Lorsqu'une opération est déclenchée depuis SAP Commerce Cloud, la plateforme prépare les informations nécessaires puis transmet une requête vers l'API correspondante du module d'intelligence artificielle. Après exécution du traitement demandé, les résultats sont retournés à SAP Commerce afin d'être exploités dans le processus métier associé.

\begin{figure}[H]
    \centering
    \includegraphics[width=0.8\textwidth]{sequence.png}
    \caption{Flux de communication entre SAP Commerce Cloud et le module IA}
    \label{fig:communication_ia}
\end{figure}

Cette organisation permet d'intégrer des capacités d'intelligence artificielle dans SAP Commerce Cloud tout en conservant une architecture modulaire, où chaque composant possède un rôle clairement défini.


\section{Extension du modèle de données SAP Commerce}

L'intégration du module d'intelligence artificielle nécessite l'adaptation du modèle de données existant de SAP Commerce afin de prendre en charge les nouvelles fonctionnalités d'enrichissement et de validation. Les extensions stockent les paramètres de traitement et les résultats générés par le module IA.

\subsection{Adaptation du modèle produit}

Le produit constitue l'élément central du processus d'enrichissement. Le modèle existant a donc été enrichi afin d'intégrer les informations nécessaires à la personnalisation des traitements d'intelligence artificielle.

Deux éléments de configuration ont été ajoutés au modèle produit. Le premier définit les informations du produit qui doivent faire l'objet d'un enrichissement ou d'une traduction. Le second précise les langues cibles dans lesquelles les contenus doivent être générés.

Ces informations permettent d'adapter le traitement selon les besoins de chaque produit. Ainsi, tous les produits ne sont pas obligés de suivre le même processus d'enrichissement, ce qui offre davantage de flexibilité dans la gestion du catalogue.

\subsection{Gestion des résultats de validation}

La validation automatique génère plusieurs informations relatives à la qualité des données produit, notamment une évaluation globale, un état de validation ainsi que les problèmes identifiés.

Afin de conserver ces résultats et d'assurer leur suivi dans le temps, un modèle spécifique a été introduit pour représenter une opération de validation. Cette séparation évite le mélange des données propres au produit avec les informations générées par le processus de contrôle qualité.

Chaque validation effectuée correspond ainsi à un résultat indépendant associé au produit concerné. Cette approche conserve l'historique des différentes analyses réalisées et suit l'évolution de la qualité des données au cours du cycle de vie du produit.

\subsection{Gestion des anomalies détectées}

Lorsqu'une validation identifie des problèmes dans les informations produit, les anomalies détectées sont conservées séparément afin de faciliter leur consultation et leur traitement.

Chaque anomalie contient les informations nécessaires à son analyse, telles que le champ concerné, le niveau de gravité, la description du problème identifié ainsi que la suggestion de correction proposée.

Cette organisation structure les résultats de validation et fournit aux utilisateurs une vision claire des améliorations nécessaires.

\subsection{Relations entre les éléments du modèle}

Les différentes extensions réalisées forment un ensemble cohérent permettant de représenter le cycle complet de traitement. Un produit peut être soumis à plusieurs validations au cours de son cycle de vie et chaque validation peut générer plusieurs anomalies.

\begin{figure}[H]
    \centering
    \includegraphics[width=0.7\textwidth]{product.png}
    \caption{Modèle de données pour la gestion de l'enrichissement et de la validation}
    \label{fig:modele_ia}
\end{figure}

Cette extension du modèle permet donc d'intégrer les capacités d'intelligence artificielle dans SAP Commerce tout en conservant une structure claire et évolutive, adaptée aux besoins d'enrichissement et de contrôle qualité des produits.





% jusq'ici , j'ai lu.

\section{Interfaces Backoffice pour l'enrichissement et la validation}

\subsection{Vue d'ensemble du Backoffice SAP Commerce}

Le Backoffice de SAP Commerce est l'interface d'administration centrale du système. C'est un espace sécurisé, accessible uniquement aux collaborateurs disposant des droits appropriés (administrateurs système, responsables de catalogues, validateurs de qualité, etc.). Contrairement à la boutique en ligne visible par les clients, le Backoffice reste un espace intérieur, pensé pour les workflows internes.

SAP Commerce fournit une interface Backoffice prête à l'emploi, appelée Backoffice Studio, dédiée à la gestion des produits, catégories, prix, stocks, commandes, etc. Cette interface repose sur le framework \texttt{ZK} \citeURL{zk_framework}, qui permet de construire des interfaces utilisateur riches en Java. Cependant, cette interface générique ne couvre pas tous les besoins métier spécifiques. Pour des processus métier particuliers, il est possible de développer des interfaces personnalisées.

Dans ce projet, deux interfaces métier ont été développées : l'une pour déclencher et suivre l'enrichissement de produits, l'autre pour consulter et agir sur les résultats de validation. Ces interfaces ont été développées en tant que composants Backoffice personnalisés, intégrés dans l'écosystème SAP Commerce.

\subsection{Interface d'enrichissement}

L'interface d'enrichissement permet aux responsables de catalogues de lancer l'enrichissement des descriptions des produits et d'en suivre l'exécution.

\subsubsection{Éléments de l'interface}

Lors du chargement de cette interface, l'utilisateur voit plusieurs éléments :

Une liste des produits sélectionnés pour enrichissement. Chaque produit est affiché avec son code, son nom, sa catégorie. L'utilisateur peut sélectionner un ou plusieurs produits via des checkboxes avant de lancer l'enrichissement. Cela permet d'enrichir des produits en masse, ce qui est beaucoup plus efficient que de traiter un produit à la fois.

Un formulaire de configuration avec deux sections principales. La première section demande à l'utilisateur quels champs enrichir. Par exemple, des checkboxes pour « description courte », « description longue », « spécifications ». L'utilisateur coche ceux qu'il souhaite enrichir. Ces informations seront stockées dans l'attribut \texttt{fieldsToTranslate} du ProductModel.

La deuxième section sert à sélectionner les langues cibles. Une liste déroulante ou un ensemble de checkboxes proposent les langues disponibles : français (FR), anglais (EN), espagnol (ES), allemand (DE), etc. L'utilisateur en sélectionne une ou plusieurs. Ces informations seront stockées dans l'attribut \texttt{targetLanguages} du ProductModel.

Un bouton « Lancer l'enrichissement » qui enclenche le processus. Quand on clique, les données du formulaire sont validées côté client (vérifier qu'au moins un champ et une langue sont sélectionnés), puis envoyées au serveur.

Une zone de statut et de progression qui affiche l'état de l'enrichissement. Avant le lancement, elle affiche « Prêt ». Après un clic sur le bouton, elle passe à « En cours ». Pendant l'enrichissement, la progression peut être affichée (par exemple, « 3 sur 10 produits traités »). Une fois terminé, elle affiche « Complété » et un résumé (nombre de produits enrichis, nombre d'erreurs).

\subsubsection{Flux de sauvegarde des données}

Quand l'utilisateur clique sur « Lancer l'enrichissement », plusieurs opérations se déroulent en arrière-plan.

Premièrement, une Facade reçoit l'appel du Backoffice. Elle valide les paramètres : les produits sélectionnés existent-ils ? L'utilisateur a-t-il le droit d'enrichir ces produits ? Les langues sélectionnées sont-elles valides ?

Deuxièmement, la Facade met à jour les ProductModels concernés, en sauvegardant la configuration : \texttt{fieldsToTranslate} reçoit la liste des champs sélectionnés, \texttt{targetLanguages} reçoit la liste des langues. Cette sauvegarde est cruciale car elle trace quelle configuration a été utilisée pour chaque produit. Si un utilisateur peut modifier la configuration après avoir lancé l'enrichissement, la traçabilité serait perdue.

Troisièmement, la Facade déclenche un processus asynchrone (un job, un workflow, ou un appel direct au service métier) qui appelle le microservice pour enrichir chaque produit. Cet asynchrone est important : l'interface Backoffice retourne immédiatement du contrôle à l'utilisateur sans attendre la fin de l'enrichissement (qui peut prendre plusieurs secondes ou minutes si le microservice est lent).

\subsubsection{Affichage des résultats}

Une fois l'enrichissement terminé, l'interface affiche les résultats. Pour chaque produit, on voit :

La description originale (par exemple, en français).

Les descriptions enrichies par le microservice, organisées par langue. Par exemple, en anglais : « High-performance running shoe designed for comfort and durability », en espagnol : « Zapato de running de alto rendimiento diseñado... », en allemand : « Hochleistungs-Laufschuh...»

L'utilisateur peut comparer la description originale avec les versions enrichies et juger de leur qualité avant validation. Si une description ne lui plaît pas, il peut la modifier manuellement (si cette fonctionnalité est implémentée) ou l'enrichissement de ce champ peut être relancé ultérieurement.

\begin{figure}[H]
	\centering
	
\includegraphics[width=0.9\textwidth]{Image1.jpg}
	\caption{Aperçu de l'interface d'enrichissement dans le Backoffice}
	\label{fig:interface_enrichissement}
\end{figure}

\subsection{Interface de validation}

L'interface de validation permet aux validateurs de qualité de consulter les résultats de validation des produits, d'identifier les anomalies et de prendre les mesures correctives appropriées.

\subsubsection{Affichage des résultats de validation}

Quand on ouvre l'interface de validation pour un produit, on voit d'abord un score de qualité global, affiché de manière saillante. Par exemple, « Score de qualité : 85/100 ». Ce score à lui seul donne une impression rapide de la qualité du produit.

Un statut synthétique accompagne le score. Par exemple : « Produit valide » (si le score est >= 90), « Produit valide avec mise en garde » (si 70 <= score < 90), « Produit à réviser » (si score < 70). Le statut est souvent codé par couleur : vert pour valide, orange pour mise en garde, rouge pour à réviser.

\subsubsection{Liste des anomalies détectées}

Sous le score, on trouve la liste des anomalies détectées lors de la validation. Chaque anomalie est affichée avec :

Le champ concerné (par exemple, « description\_en »).

Le niveau de sévérité (INFO, AVERTISSEMENT, ERREUR). Une anomalie de sévérité ERROR bloque possiblement certaines actions (publication du produit) ; une AVERTISSEMENT est signalée mais moins critique ; une INFO est informative.

Une description du problème en termes lisibles (par exemple, « La description en anglais est inférieure à 50 caractères »).

Une suggestion de correction (par exemple, « Ajouter au moins 30 caractères pour une description complète »).

\subsubsection{Actions correctives}

Pour chaque anomalie, l'interface propose une action. Dans les cas simples, il est possible d'accepter une suggestion de correction automatique, ce qui applique la modification au produit immédiatement. Par exemple, si le microservice suggère « Ajouter un point à la fin de la description », un clic sur « Appliquer la correction » effectue cette modification.

Pour les cas plus complexes, l'utilisateur peut éditer manuellement le champ concerné via un éditeur de texte intégré à l'interface.

\subsubsection{Workflow de validation}

Le processus complet ressemble à ceci. L'utilisateur ouvre l'interface de validation et cherche un produit qui nécessite validation. Il sélectionne le produit. Le système récupère le ProductValidationResult le plus récent associé à ce produit et affiche le score, le statut et les anomalies.

L'utilisateur lit les anomalies et décide d'agir. Si les anomalies sont mineures et que les suggestions automatiques conviennent, il clique sur « Appliquer toutes les corrections ». Le système applique les corrections et marque le produit comme « Validé ».

Si l'utilisateur n'est pas d'accord avec une suggestion, il peut éditer le champ concerné manuellement, puis relancer la validation pour voir si les corrections éliminent l'anomalie.

Si un score est critique (< 50, par exemple), l'interface peut empêcher certaines actions comme la publication du produit jusqu'à correction.

\begin{figure}[H]
	\centering
	
\includegraphics[width=0.9\textwidth]{Image2.jpg}
	\caption{Aperçu de l'interface de validation dans le Backoffice}
	\label{fig:interface_validation}
\end{figure}

Les deux interfaces d'enrichissement et de validation permettent à l'utilisateur de piloter manuellement les processus. Cependant, pour une véritable automatisation et pour éviter l'intervention humaine répétitive, SAP Commerce dispose d'un moteur de workflows. C'est le sujet de la section suivante.

\section{Workflows métier et automatisation}

\subsection{Workflows dans SAP Commerce}

Un workflow est une séquence d'étapes automatisées qui exécutent un processus métier selon une logique prédéfinie. Dans le contexte de SAP Commerce, le moteur de workflow gère l'ordonnancement des étapes, les transitions conditionnelles, l'appel des actions métier et l'enregistrement de l'historique.

Pourquoi utiliser un workflow au lieu d'écrire simplement une série d'appels de services ? Plusieurs avantages le justifient. D'abord, un workflow rend le processus visible et compréhensible pour les non-développeurs : un responsable métier peut lire la définition du workflow et comprendre comment fonctionnent les processus. Ensuite, un workflow est flexible : modifier le processus ne nécessite pas de redéployer du code Java ; on modifie la configuration du workflow. Troisièmement, les workflows offrent des capacités avancées : parallélisation d'étapes, gestion des erreurs, relance automatique, etc. Enfin, les workflows fournissent une traçabilité complète : chaque étape est loggée, et on peut consulter l'historique d'exécution.

SAP Commerce utilise un moteur de workflow déclaratif. Les workflows ne sont pas écrits en Java, mais définis via un format de configuration (XML ou une interface graphique appelée Workflow Builder). On décrit : « L'étape A s'exécute, puis on teste une condition, si elle est vraie on va à l'étape B, sinon à l'étape C ». Le moteur lit cette configuration et l'exécute.

\subsubsection{Implémentation dans le projet}

Dans ce projet, les workflows ont été définis via le mécanisme d'impex de SAP Commerce. Impex est un format texte de configuration qui permet de définir des données, des configurations, et des workflows. Les définitions de workflow ont été rédigées en impex et importées dans SAP Commerce lors du déploiement.

Chaque étape du workflow fait référence à une action métier. Une action métier est une classe Java implémentant une interface standard (par exemple, \texttt{WorkflowAction}). Quand le workflow atteint cette étape, il instancie la classe et exécute sa méthode \texttt{execute()}.

\subsection{Workflow d'enrichissement}

Le workflow d'enrichissement automatise complètement le processus d'enrichissement d'un produit, du déclenchement jusqu'à la persistance des résultats.

\subsubsection{Étapes du workflow}

Le workflow se déroule selon une séquence bien définie :

\textbf{Étape 1 : Déclenchement.} Le workflow est déclenché quand l'utilisateur clique sur « Lancer l'enrichissement » dans l'interface Backoffice. À ce moment, les ProductModels concernés ont déjà été mis à jour avec la configuration (champs à enrichir, langues cibles). Le workflow reçoit en paramètre la liste des IDs de produits à enrichir.

\textbf{Étape 2 : Préparation des données.} Une action métier appelée \texttt{ProductEnrichmentAutomatedAction} s'exécute. Elle récupère le ProductModel depuis la base de données, extrait les informations pertinentes (code, nom, descriptions actuelles, champs à traduire, langues cibles) et construit un objet \texttt{ProductEnrichmentRequest} prêt à envoyer au microservice.

\textbf{Étape 3 : Appel au microservice.} Une seconde action (ou la même, selon l'implémentation) utilise RestTemplate pour envoyer la requête au microservice et récupère la réponse JSON, désérialisée en objet \texttt{ProductEnrichmentResponse} par Jackson.

\textbf{Étape 4 : Traitement de la réponse.} L'action vérifie que la réponse indique un succès (\texttt{success = true}). Si oui, elle procède à l'étape suivante. Si non (par exemple, si le microservice a signalé une erreur métier), elle redirige le workflow vers une étape de gestion d'erreur.

\textbf{Étape 5 : Mise à jour du produit.} L'action extrait les descriptions enrichies de la réponse (descriptions par langue) et les sauvegarde dans le ProductModel. Concrètement, cela signifie mettre à jour les attributs \texttt{description\_en}, \texttt{description\_es}, etc. avec les valeurs retournées par le microservice. Ces modifications sont persistées en base de données via le ModelService.

\textbf{Étape 6 : Déclenchement du workflow de validation.} Après enrichissement réussi, le workflow d'enrichissement déclenche automatiquement le workflow de validation pour le même produit (voir section 4.6.3). Cela crée une chaîne d'automatisation : enrichir puis valider immédiatement.

\textbf{Étape 7 : Fin du workflow.} Une étape de conclusion enregistre que le workflow s'est terminé avec succès, log un message et éventuellement notifie l'utilisateur que l'enrichissement est terminé (voir section 4.8).

\subsubsection{Gestion des erreurs}

À plusieurs points du workflow, des erreurs peuvent survenir. Le microservice peut être inaccessible (serveur web down, timeout réseau). La réponse du microservice peut être malformée (JSON invalide, champs manquants). La base de données peut être inaccessible lors de la sauvegarde.

Pour chacune de ces situations, le workflow contient une branche de gestion d'erreur. Typiquement, si un appel au microservice échoue, le workflow saute à une étape « Gestion d'erreur » qui :

Enregistre l'erreur dans les logs du système.

Notifie l'administrateur système qu'un enrichissement a échoué et a besoin d'intervention.

Marque le produit avec un statut « Enrichissement échoué - Nouvelle tentative requise ».

Optionnellement, déclenche une tentative automatique après un délai (par exemple, réessayer après 5 minutes).

Sans cette gestion robuste des erreurs, un problème réseau passager bloquerait indéfiniment l'enrichissement.

\begin{figure}[H]
	\centering
	\includegraphics[width=0.8\textwidth]{activités.png}
	\caption{Diagramme d'activités du workflow d'enrichissement}
	\label{fig:workflow_activites}
\end{figure}

\subsection{Workflow de validation}

Le workflow de validation fonctionne de manière similaire au workflow d'enrichissement, avec des différences métier importantes.

\subsubsection{Étapes du workflow}

\textbf{Étape 1 : Déclenchement.} Le workflow de validation est déclenché soit manuellement par l'utilisateur (via l'interface de validation), soit automatiquement à la fin du workflow d'enrichissement. Il reçoit en paramètre l'ID du produit à valider.

\textbf{Étape 2 : Appel au microservice de validation.} Une action appelle l'endpoint de validation du microservice avec un objet \texttt{ProductValidationRequest} contenant les descriptions actuelles du produit (en toutes les langues) et d'autres métadonnées.

\textbf{Étape 3 : Réception du score et des anomalies.} Le microservice retourne une réponse contenant un score de qualité (0-100), un statut et une liste d'anomalies détectées.

\textbf{Étape 4 : Persistance des résultats.} Une action crée un nouvel enregistrement ProductValidationResult avec le score et le statut reçus. Pour chaque anomalie retournée par le microservice, une instance de ValidationAnomaly est créée et liée au ProductValidationResult. Ces entités sont sauvegardées en base de données.

\textbf{Étape 5 : Prise de décision basée sur le score.} Le workflow teste le score. Si le score est >= 80 (par exemple), le produit est marqué comme « Validé ». Si 60 <= score < 80, le statut devient « Avertissement ». Si score < 60, le statut devient « À réviser ». Ce seuil peut être configuré.

\textbf{Étape 6 : Notification.} Un email est généré et envoyé au responsable de la validation (voir section 4.8) avec un résumé : score, statut, anomalies critiques. Si le statut est « À réviser », un email additionnel peut être envoyé au responsable produit pour action.

\textbf{Étape 7 : Fin du workflow.} Enregistrement de la completion du workflow.

\subsubsection{Différences avec le workflow d'enrichissement}

Contrairement au workflow d'enrichissement qui modifie le ProductModel, le workflow de validation ne modifie pas les descriptions : il les analyse seulement. Le workflow d'enrichissement envoie des données au microservice et en récupère d'autres ; le workflow de validation envoie des données et reçoit une évaluation, pas des données de remplacement.

De plus, le workflow de validation génère toujours une notification, tandis que le workflow d'enrichissement peut la générer optionnellement.

Maintenant que les workflows orchestrent l'automation, il est instructif de voir comment tout s'imbrique. La section suivante trace le flux complet d'une requête utilisateur à travers l'ensemble du système.

\section{Flux complet d'exécution et orchestration}

\subsection{Scénario concret : enrichissement d'un produit}

Pour bien comprendre comment tous les éléments décrits précédemment s'articulent, traçons le parcours complet d'une requête utilisateur.

Une responsable de catalogue, nommée Marie, travaille dans le Backoffice. Elle gère un produit phare de l'entreprise, un article de chaussure sportive avec le code « CHAUSSURE-001 » et le nom « Chaussure de sport Running ». Ce produit a une description courte en français : « Chaussure de sport de haute qualité ». Marie souhaite enrichir ce produit en générant des descriptions de qualité professionnelle en anglais, espagnol et allemand.

Elle navigue jusqu'à l'interface d'enrichissement. Elle sélectionne le produit « CHAUSSURE-001 ». Elle coche le champ « Description courte » pour enrichissement. Elle sélectionne les langues « Anglais », « Espagnol », « Allemand ». Elle clique sur le bouton « Lancer l'enrichissement ».

\subsection{Étapes techniques du flux complet}

\subsubsection{Étape 1 : Validation et préparation (Couche Présentation/Facade)}

L'interface Backoffice envoie une requête HTTP vers le serveur SAP Commerce, transmettant les paramètres : productCode = « CHAUSSURE-001 », fieldsToTranslate = [« description »], targetLanguages = [« en », « es », « de »].

Une Facade appellée \texttt{ProductEnrichmentFacade} reçoit cette requête. Elle effectue des validations préliminaires :

Existe-t-il un produit avec le code CHAUSSURE-001 ? (Oui)

L'utilisateur actuel (Marie) a-t-elle les droits pour enrichir les produits ? (Oui)

Les langues « en », « es », « de » sont-elles des codes valides ? (Oui)

Si une de ces validations échoue, la Facade retourne immédiatement une erreur à l'interface, qui l'affiche à l'utilisateur.

\subsubsection{Étape 2 : Sauvegarde de la configuration}

Après validation, la Facade met à jour le ProductModel de CHAUSSURE-001 en base de données :

\texttt{fieldsToTranslate} = [« description »]

\texttt{targetLanguages} = [« en », « es », « de »]

Cette sauvegarde est persistée via le ModelService, qui utilise un DAO interne pour mettre à jour la base de données.

\subsubsection{Étape 3 : Déclenchement du workflow d'enrichissement}

Après cette sauvegarde, la Facade déclenche le workflow d'enrichissement, en lui passant l'ID du produit CHAUSSURE-001.

\subsubsection{Étape 4 : Orchestration métier (Couche Service)}

Le workflow d'enrichissement s'exécute. La première action métier (\texttt{ProductEnrichmentAutomatedAction}) s'exécute. Elle :

Récupère le ProductModel CHAUSSURE-001 de la base de données.

Extrait les informations :
\begin{itemize}
	\item productCode = « CHAUSSURE-001 »
	\item productName = « Chaussure de sport Running »
	\item description = « Chaussure de sport de haute qualité »
	\item fieldsToTranslate = [« description »]
	\item targetLanguages = [« en », « es », « de »]
\end{itemize}

Construit un objet DTO de type \texttt{ProductEnrichmentRequest} avec ces données.

\subsubsection{Étape 5 : Sérialisation et communication HTTP}

L'action appelle un service métier \texttt{ProductEnrichmentService}, en lui transmettant le ProductModel.

Le service utilise \texttt{RestTemplate} \citeURL{spring_framework} pour appeler le microservice. \texttt{RestTemplate} :

Prend l'objet \texttt{ProductEnrichmentRequest} et utilise la bibliothèque \texttt{Jackson} \citeURL{jackson_repo} pour le sérialiser en JSON :

\begin{verbatim}
{
  "productCode": "CHAUSSURE-001",
  "productName": "Chaussure de sport Running",
  "description": "Chaussure de sport de haute qualité",
  "fieldsToTranslate": ["description"],
  "targetLanguages": ["en", "es", "de"]
}
\end{verbatim}

Construit une requête HTTP POST vers l'URL du microservice (par exemple, \texttt{http://microservice-ia:8000/api/v1/enrich}).

Ajoute les en-têtes HTTP appropriés : \texttt{Content-Type: application/json}, \texttt{Authorization: Bearer <token>}, etc.

Établit une connexion TCP vers le serveur microservice.

Envoie la requête avec le JSON en corps.

Attend la réponse pendant au maximum 30 secondes (timeout configuré).

\subsubsection{Étape 6 : Traitement du microservice}

Le microservice Python reçoit la requête. Il exécute le traitement décrit au chapitre 3 : génération de descriptions enrichies en utilisant les modèles de langage, traduction en les trois langues cibles. Après environ 2-3 secondes, il construit une réponse JSON :

\begin{verbatim}
{
  "success": true,
  "enrichedData": {
    "en": "High-performance sports shoe combining comfort, durability, 
           and style. Perfect for everyday running and athletic activities.",
    "es": "Zapato deportivo de alto rendimiento que combina comodidad, 
           durabilidad y estilo. Perfecto para correr diariamente...",
    "de": "Hochleistungs-Sportschuh, der Komfort, Haltbarkeit und Stil 
          kombiniert. Perfekt für alltägliches Laufen..."
  },
  "timestamp": "2024-01-15T14:30:45Z",
  "processingTimeMs": 2850
}
\end{verbatim}

Il envoie cette réponse en tant que corps d'une réponse HTTP 200 OK.

\subsubsection{Étape 7 : Désérialisation de la réponse}

RestTemplate reçoit la réponse HTTP. Jackson désérialise le JSON en un objet Java de type \texttt{ProductEnrichmentResponse} :

\begin{verbatim}
ProductEnrichmentResponse response = new ProductEnrichmentResponse();
response.setSuccess(true);
Map<String, String> enrichedData = new HashMap<>();
enrichedData.put("en", "High-performance sports shoe...");
enrichedData.put("es", "Zapato deportivo...");
enrichedData.put("de", "Hochleistungs-Sportschuh...");
response.setEnrichedData(enrichedData);
response.setTimestamp(LocalDateTime.parse("2024-01-15T14:30:45Z"));
\end{verbatim}

Le service métier reçoit cet objet Java structuré.

\subsubsection{Étape 8 : Mise à jour du ProductModel}

Le service métier extrait les descriptions enrichies de la réponse. Il met à jour le ProductModel CHAUSSURE-001 :

\texttt{description\_en} = « High-performance sports shoe... »

\texttt{description\_es} = « Zapato deportivo... »

\texttt{description\_de} = « Hochleistungs-Sportschuh... »

Le ModelService persiste ces modifications en base de données via les DAOs.

\subsubsection{Étape 9 : Déclenchement du workflow de validation}

À la fin du workflow d'enrichissement, le workflow de validation est déclenché automatiquement pour le même produit CHAUSSURE-001.

\subsubsection{Étape 10 : Validation des descriptions}

Le workflow de validation appelle le microservice avec l'endpoint de validation, transmettant les descriptions enrichies. Le microservice analyse chaque description selon ses critères de qualité (longueur, richesse du vocabulaire, cohérence, complétude) et retourne :

\begin{verbatim}
{
  "score": 87,
  "status": "VALID_WITH_WARNINGS",
  "anomalies": [
    {
      "field": "description_es",
      "severity": "WARNING",
      "message": "Description espagnole un peu courte (65 caractères)",
      "suggestion": "Ajouter des détails sur les matériaux ou les performances"
    }
  ],
  "timestamp": "2024-01-15T14:31:00Z"
}
\end{verbatim}

\subsubsection{Étape 11 : Persistance des résultats de validation}

Le workflow de validation reçoit cette réponse. Il crée un nouvel enregistrement ProductValidationResult :

ID = 42 (généré automatiquement)

Product = CHAUSSURE-001

Score = 87

Status = VALID\_WITH\_WARNINGS

ValidationDate = 2024-01-15 14:31:00

Il crée aussi une entité ValidationAnomaly :

ID = 101

ProductValidationResult = 42

Field = description\_es

Severity = WARNING

Message = « Description espagnole un peu courte (65 caractères) »

Suggestion = « Ajouter des détails sur les matériaux ou les performances »

Ces deux entités sont sauvegardées en base de données.

\subsubsection{Étape 12 : Notification par email}

Après validation complète, le workflow déclenche l'envoi d'un email (voir section 4.8). Un template Velocity génère le contenu :

\begin{verbatim}
Objet : Enrichissement et validation du produit CHAUSSURE-001

Bonjour,

Le produit "Chaussure de sport Running" (CHAUSSURE-001) a été 
enrichi et validé avec succès.

Score de qualité : 87/100
Statut : Valide avec mise en garde

Anomalies détectées :
- description_es (AVERTISSEMENT) : Description espagnole un peu courte. 
  Suggestion : Ajouter des détails sur les matériaux ou les performances.

Vous pouvez consulter les détails dans l'interface de validation.

Cordialement,
Système d'enrichissement de produits
\end{verbatim}

Cet email est envoyé à la responsable de validation (ou à Marie si elle l'a demandé).

\subsubsection{Étape 13 : Retour vers l'interface}

L'interface Backoffice, qui attendait le résultat de la Facade (depuis l'étape 1), reçoit enfin une réponse. L'opération asynchrone en arrière-plan s'est terminée et le résultat remonte vers l'interface.

L'interface affiche :

« Enrichissement terminé avec succès »

Les descriptions enrichies en trois langues

Le score de validation (87/100)

Les anomalies détectées (avec possibilité d'appliquer les corrections suggérées)

Marie peut maintenant consulter les descriptions enrichies. Elle voit que la description en espagnol est un peu courte (selon l'avis du microservice), mais dans l'ensemble le résultat est satisfaisant. Elle clique sur « Approuver et publier » pour marquer le produit comme validé et le rendre visible sur la boutique en ligne.

\begin{figure}[h]
	\centering
	\includegraphics[width=\textwidth]{sequence}
	\caption{Diagramme de séquence global de l'interaction utilisateur et des processus système}
	\label{fig:sequence_globale_processus}
\end{figure}

Cet enchaînement illustre l'interopérabilité entre les composants du système : l'interface, les Facades, les services métier, la sérialisation des données, la communication HTTP, le traitement du microservice, la persistance, les workflows et les notifications. Aucun de ces éléments ne fonctionne de manière isolée ; ils forment un ensemble cohérent où chaque maillon est indispensable.

\section{Système de notification par courrier électronique}

\subsection{Moteur d'email dans SAP Commerce}

SAP Commerce intègre un moteur d'email complet qui gère l'envoi de courriers électroniques à partir du système. Ce moteur découple la génération du contenu de l'email de son envoi via SMTP, ce qui offre plusieurs avantages en termes de performance et de fiabilité.

La logique est la suivante : quand un événement métier survient (fin d'enrichissement, fin de validation, erreur), un processus génère le contenu de l'email (le corps du message, le sujet, les destinataires). Cet email est ajouté à une queue (une file d'attente). Un worker en arrière-plan consomme cette queue et envoie les emails via le serveur SMTP configuré.

Cette séparation présente plusieurs avantages. D'abord, la génération du contenu peut être lente (si elle implique des requêtes en base de données, des appels externes, etc.), mais cela n'impacte pas l'interface utilisateur : l'interface reçoit le contrôle immédiatement, tandis que la génération du contenu s'effectue en arrière-plan. Ensuite, si le serveur SMTP est temporairement inaccessible, les emails restent dans la queue et sont renvoyés ultérieurement une fois que le serveur SMTP redevient disponible. Aucun email n'est perdu.

\subsection{Moteurs de template Velocity}

Plutôt que de générer les emails via des concaténations de chaînes de caractères en Java (ce qui serait verbeux et peu flexible), SAP Commerce utilise le moteur \texttt{Velocity} \citeURL{velocity_engine}.

\texttt{Velocity} \citeURL{velocity_engine} est un moteur de template simple mais puissant. Un template Velocity est un fichier texte contenant du texte fixe et des directives de template délimitées par des symboles spéciaux (généralement \texttt{\$} pour les variables, \texttt{\#} pour les directives). Par exemple :

\begin{verbatim}
Objet : Enrichissement du produit $productName

Bonjour,

Le produit "$productName" (code : $productCode) a été enrichi 
avec succès.

Descriptions générées :
#foreach($language in $targetLanguages)
  - $language : ${enrichedDescriptions.get($language)}
#end

Score de validation : $validationScore/100

#if($validationScore >= 80)
Produit validé. Prêt pour publication.
#else
Avertissement : le produit nécessite une révision avant publication.
#end

Anomalies détectées :
#foreach($anomaly in $anomalies)
  - [$anomaly.field] ($anomaly.severity) : $anomaly.message
#end

Cordialement,
L'équipe de gestion des produits
\end{verbatim}

Quand on fournit des données à ce template (productName = « Chaussure de sport Running », productCode = « CHAUSSURE-001 », targetLanguages = [« en », « es », « de »], validationScore = 87, etc.), Velocity remplace les variables par leurs valeurs, évalue les directives conditionnelles (\texttt{\#if}), itère sur les listes (\texttt{\#foreach}), et génère le contenu final.

Pourquoi utiliser Velocity plutôt que de coder l'email directement en Java ? Plusieurs raisons. D'abord, les templates Velocity sont écrits en texte brut, facilement éditables. Un responsable métier peut modifier un template sans compétences Java. Ensuite, modifier un template ne nécessite pas une recompilation et un redéploiement de l'application : SAP Commerce peut recharger les templates dynamiquement. Enfin, les templates permettent une séparation nette entre la présentation (ce que dit l'email) et la logique (ce qui construit les données).

\subsection{Configuration SMTP et Mailtrap}

SMTP (Simple Mail Transfer Protocol) est le protocole standard pour envoyer des emails. Pour envoyer un email, SAP Commerce a besoin de la configuration d'un serveur SMTP : l'adresse du serveur, le port, les identifiants (login et mot de passe), etc.

En production, le serveur SMTP est généralement un serveur mail de l'entreprise (Google Workspace, Microsoft Exchange, ou un serveur mail auto-hébergé). Cependant, en environnement de développement et de test, envoyer de vrais emails est problématique : on risque d'envoyer des emails de test à des adresses email réelles d'utilisateurs.

Pour résoudre ce problème, des services comme Mailtrap existent. Mailtrap simule un serveur SMTP mais, au lieu de réellement envoyer les emails, les capture et les stocke pour inspection. Cela permet aux développeurs de tester le fonctionnement du moteur d'email sans risque d'envoyer des emails indésirables.

Dans ce projet, Mailtrap a été utilisé pour les environnements de développement et de test. La configuration SAP Commerce pointait vers le serveur Mailtrap. Lors du lancement d'un enrichissement ou d'une validation en développement, les emails générés étaient capturés par Mailtrap et visibles via son interface web. En production, la configuration serait changée pour pointer vers le vrai serveur SMTP de l'entreprise.

\subsection{Processus complet de notification}

\subsubsection{Déclenchement}

Quand le workflow de validation se termine, une action métier appelle le moteur d'email de SAP Commerce. Elle spécifie :

Le template à utiliser (par exemple, \texttt{enrichmentAndValidationNotification})

Les données à passer au template (ProductModel, ProductValidationResult, ValidationAnomalies, etc.)

Le destinataire de l'email (par exemple, l'adresse email de la responsable de validation)

Le sujet de l'email (optionnel si le template le définit)

\subsubsection{Génération du contenu}

SAP Commerce charge le template spécifié (du fichier système ou de la base de données). Il exécute le moteur Velocity avec les données fournies. Velocity remplace les variables par leurs valeurs et évalue les directives.

Par exemple, si le template contient :

\begin{verbatim}
#foreach($anomaly in $anomalies)
  - $anomaly.field : $anomaly.message
#end
\end{verbatim}

Et que \texttt{\$anomalies} contient deux anomalies, Velocity génère :

\begin{verbatim}
  - description_es : Description espagnole un peu courte
  - description_en : Point manquant à la fin
\end{verbatim}

Le contenu final généré ressemble à un email normal, avec un sujet, un corps, et éventuellement des pièces jointes (PDF, images, etc.).

\subsubsection{Ajout à la queue et envoi}

L'email généré est créé en tant qu'objet (sujet, corps, destinataires, pièces jointes) et ajouté à une queue de traitement. Un worker (processus en arrière-plan ou job batch) consomme cette queue périodiquement.

Pour chaque email dans la queue, le worker utilise la configuration SMTP pour ouvrir une connexion au serveur SMTP, authentifie SAP Commerce, et envoie l'email. Si l'envoi échoue (serveur inaccessible, destinataire invalide), le worker enregistre l'erreur et peut relancer l'envoi ultérieurement.

\subsubsection{Cas d'usage dans le projet}

Trois types de notifications ont été implémentées :

\textbf{Email après enrichissement réussi :} « Votre enrichissement de produit est terminé. 3 produits enrichis avec succès. 0 erreurs. »

\textbf{Email après validation complète :} « Validation terminée. Produit CHAUSSURE-001 : Score 87/100. 1 avertissement détecté. Voir les détails dans le Backoffice. »

\textbf{Email d'alerte en cas d'erreur :} « Enrichissement du produit CHAUSSURE-001 a échoué. Le microservice n'a pas répondu dans le délai imparti. Veuillez réessayer ou contacter l'administrateur système. »

\begin{figure}[h]
	\centering
	\fbox{\begin{minipage}{0.9\textwidth}
			\textbf{[Figure 4.9 : Diagramme de flux - Système de notifications]}
			
			Insérer un diagramme montrant :
			\begin{itemize}
				\item Workflow ou Service métier
				\item Appel au moteur Email avec template et données
				\item Moteur Velocity génère contenu
				\item Email créé (sujet, corps, destinataires)
				\item Ajout à la queue d'emails
				\item Worker consomme queue
				\item Configuration SMTP
				\item Connexion au serveur SMTP
				\item Envoi effectif de l'email
				\item Utilisateur reçoit email
				\item Feedback de succès/erreur
			\end{itemize}
		\end{minipage}}
\end{figure}

\section{Conclusion du chapitre}

L'intégration dépasse le simple appel HTTP pour embrasser une dimension systémique via six axes complémentaires.

Ce chapitre a présenté ces six axes d'intégration :

Premièrement, l'architecture logicielle de SAP Commerce elle-même, organisée en couches (Présentation, Facade, Service, DAO, Modèle), où chaque couche a une responsabilité bien définie. Les développements du projet se situent dans toutes les couches.

Deuxièmement, la communication technique entre SAP Commerce et le microservice via REST, RestTemplate et Jackson. Cette communication est robuste, gère les erreurs, et persiste les résultats.

Troisièmement, l'extension du modèle de données pour accueillir les nouvelles entités métier : l'enrichissement de ProductModel avec des attributs de configuration, et la création de ProductValidationResult et ValidationAnomaly pour persister les résultats de validation avec historique.

Quatrièmement, les interfaces Backoffice personnalisées qui permettent aux utilisateurs métier de déclencher l'enrichissement et de consulter les résultats de validation. Ces interfaces rendent accessibles les capacités du microservice sans nécessiter de connaissances techniques.

Cinquièmement, l'automatisation via les workflows métier. Les workflows d'enrichissement et de validation orchestrent les étapes, gèrent les erreurs, et assurent la traçabilité. Ils rendent le processus visible et modifiable sans intervention de développeurs.

Sixièmement, le système de notifications par email, utilisant le moteur Email de SAP Commerce, les templates Velocity, et la configuration SMTP. Les notifications informent les utilisateurs en temps réel de la progression et des anomalies.

Au-delà de ces six axes, ce qui rend cette intégration complète c'est leur cohésion. Aucun de ces éléments ne fonctionne de manière isolée. L'interface déclenche un workflow, qui appelle un service, qui appelle le microservice via RestTemplate, qui est désérialisé par Jackson, qui met à jour le modèle via un DAO, qui persiste les résultats, qui sont ensuite affichés par l'interface, et dont on notifie l'utilisateur par email. C'est cette orchestration globale qui constitue une véritable intégration.

Le chapitre suivant aborde la validation de cette intégration : comment s'assurer que le système fonctionne correctement, que les enrichissements sont de qualité, que les erreurs sont gérées robustement, et que les performances sont acceptables. Cela passe par les stratégies de test, les expérimentations, et l'analyse des résultats.

